\section{Evaluation}

% Lead owner: Keerthi
% Exp 3 write-up: Sreeja
% Exp 5 interpretation: shared

Frame the evaluation around one rule: each experiment must answer one explicit
question. The goal is not to enumerate every chart, but to show whether PBCP
solves the governance-timing problem and whether its intent-aware signals add
value beyond simpler baselines.

\subsection{Experimental Setup}

Summarize the controlled benchmark briefly: 500 workloads, 28,423 runs, six
workload classes, historical run traces, and the intervention stages evaluated by
the system. Keep this concise and reserve schema detail for an appendix or
artifact note if needed.

\subsection{Exp 0: Can Simulation Predict Utilization Accurately?}

Report calibration quality and the corrected error metrics that justify using the
simulator as a pre-execution signal. This subsection should anchor later claims
by showing that predicted utilization and cost are close enough to support
intervention decisions.

\subsection{Exp 1: Can Pre-Provision Intervention Reduce Waste?}

Use this subsection to compare full PBCP against the baseline methods and to
highlight the 20-node ETL showcase scenario. Keep the question narrow: does
pre-billing intervention prevent waste before resources are launched?

\subsection{Exp 2: Can Runtime Governance Catch Hidden Failures?}

Describe the scenarios where pre-execution reasoning is not sufficient and
runtime correction still matters. Focus on what kinds of hidden failures were
caught and how much cost was prevented after launch.

\subsection{Exp 3: Does IFS Outperform CPU-Threshold Detection?}

This subsection is primarily Sreeja's write-up. Compare IFS-based anomaly
detection against the utilization-only baseline, report precision/recall/F1, and
discuss the mismatch-subgroup near-fail honestly rather than hiding it. The
current result should be framed precisely: the primary detection-quality gate
passes clearly, while the subgroup comparison is mixed and should be interpreted
as a limitation rather than ignored.

% Owner: Sreeja

This insert is reserved for the Exp 3 anomaly-detection write-up. It should
compare IFS against the CPU-threshold baseline directly, report the F1
improvement clearly, and discuss the subgroup nuance as a limitation rather than
over-claiming generality.


\subsection{Exp 5: What Is Overall System Effectiveness?}

Interpret the system roll-up using CPS, ESR, and IFS together. This subsection
should explain how the full system performs across workload types while keeping
the metric relationship clear: prevention, execution safety, and behavioral
alignment.

\subsection{Exp 6: Does Policy Learning Improve Over Time?}

Use the convergence experiment to show whether the learning loop improves beyond
frozen or partial-system baselines. The emphasis should be on longitudinal
improvement, not on UI presentation of the curve.
